\chapter*{Introduction Générale}
\addcontentsline{toc}{chapter}{Introduction Générale}

À l'ère de la transformation numérique, la modernisation des outils de gestion constitue un levier de performance incontournable pour les grandes structures industrielles. Les systèmes d'information traditionnels, souvent basés sur des processus manuels ou des outils décentralisés, montrent aujourd'hui leurs limites face aux exigences de réactivité, de traçabilité et de sécurité des organisations modernes. La gestion des actifs physiques et du patrimoine matériel représente, à ce titre, un enjeu stratégique majeur pour optimiser l'exploitation des ressources et garantir un suivi rigoureux des investissements.

C'est dans ce contexte que s'inscrit le présent projet, réalisé au sein de la Société Tunisienne de l'Électricité et du Gaz (STEG). En tant qu'acteur national pivot du secteur énergétique, la société s'appuie sur une infrastructure technique et immobilière d'envergure, répartie sur de nombreuses directions régionales et locaux de stockage. L'absence d'une solution logicielle unifiée et accessible en temps réel pour le suivi des immobilisations, telles que les parcs de véhicules, les équipements électriques ou les structures physiques, engendre des difficultés d'inventaire et complexifie la coordination entre les différents services responsables.

L'objectif principal de ce travail consiste donc à concevoir et à développer une application web dédiée à la gestion automatisée des immobilisations. Cette solution doit permettre de recenser de manière centralisée le patrimoine de la structure, de suivre avec précision l'état et la localisation de chaque type de matériel, et de fluidifier le processus de contrôle des inventaires. Pour répondre aux besoins de sécurité et de répartition des tâches au sein de l'entreprise, la plateforme doit également proposer un contrôle d'accès strict basé sur différents profils d'utilisateurs.

Pour rendre compte de la réalisation de ce système, le présent rapport s'organise en plusieurs chapitres complémentaires. La première partie expose le contexte général du projet et présente la structure d'accueil ainsi que la problématique identifiée. La deuxième partie est consacrée à l'analyse détaillée des exigences fonctionnelles et non fonctionnelles, tout en spécifiant les profils des acteurs et l'écosystème logiciel retenu. Enfin, les dernières sections détaillent la phase de conception technique à travers la modélisation des bases de données et des processus, avant de présenter l'implémentation logicielle, les interfaces de l'application et les résultats des tests d'intégration.